% Sections 4 and 5 of zkMetal paper
% USENIX Security format (2-column, 10pt)

\section{GPU Scheduling Pathology}
\label{sec:scheduling}

We report a previously undocumented performance pathology on Apple M-series GPUs
where specific threadgroup grid dimensions cause 10--30$\times$ latency
regressions in compute dispatches. This phenomenon was discovered independently
across three elliptic curves during MSM optimization and is not documented in
Apple's Metal Best Practices Guide~\cite{apple-metal-best-practices} or any
prior literature we are aware of.

\subsection{Phenomenon}

MSM bucket accumulation dispatches one GPU thread per bucket, where the bucket
count is determined by the Pippenger window width parameter $w$:
$\mathit{nBuckets} = 2^{w-1}$ (with signed-digit decomposition) or $2^w - 1$
(unsigned). We observed that specific $(w, \text{curve})$ combinations produce
catastrophic slowdowns far exceeding what register pressure or occupancy models
predict.

\paragraph{BN254 ($w{=}14$).}
At window width 14 with unsigned decomposition, $\mathit{nBuckets} = 16{,}384 =
2^{14}$. MSM latency jumps to 7{,}000\,ms---a 140$\times$ regression over the
expected ${\sim}50$\,ms based on extrapolation from adjacent window sizes. This
was the first observed instance and initially attributed to a correctness bug
before being confirmed as a pure performance anomaly.

\paragraph{secp256k1 ($w{=}12$).}
With signed-digit decomposition, $w{=}12$ yields $\mathit{nBuckets} = 2{,}049$.
MSM at $2^{14}$ points takes 78\,ms versus 21\,ms at $w{=}13$ ($\mathit{nBuckets} =
4{,}097$), a 3.7$\times$ regression. Similarly, $w{=}11$ ($\mathit{nBuckets} = 1{,}025$)
and $w{=}15$ ($\mathit{nBuckets} = 16{,}385$) are both pathological (131\,ms and
244\,ms respectively at $2^{14}$).

\paragraph{BLS12-377 ($w{=}14$).}
This curve uses a 377-bit base field requiring 12$\times$32-bit limbs per field
element, so each GPU thread consumes substantially more registers. At $w{=}14$,
MSM on $2^{17}$ points takes 17{,}000\,ms versus 569\,ms at $w{=}15$---a
30$\times$ regression.

\paragraph{GLV interaction.}
The GLV endomorphism~\cite{gallant-lambert-vanstone} doubles the effective point
count, which can push the bucket count into a pathological band. On secp256k1
with GLV enabled, the effective $2^{18}$ points at $w{=}16$ yield 32{,}768
buckets, causing 13{,}569\,ms versus 1{,}133\,ms without GLV---a 12$\times$
regression that entirely negates GLV's theoretical 22\% benefit. This forced us
to disable GLV on secp256k1 by default.

\begin{table}[t]
\centering
\small
\begin{tabular}{lrrrl}
\toprule
\textbf{Curve} & $w$ & \textbf{Buckets} & \textbf{Regression} & \textbf{Trigger} \\
\midrule
BN254     & 14 & 16{,}384  & 140$\times$ & $2^{14}$ \\
secp256k1 & 11 & 1{,}025   & 6.2$\times$ & $\approx 2^{10}{+}1$ \\
secp256k1 & 12 & 2{,}049   & 3.7$\times$ & $\approx 2^{11}{+}1$ \\
secp256k1 & 15 & 16{,}385  & 11.6$\times$ & $\approx 2^{14}{+}1$ \\
secp256k1+GLV & 16 & 32{,}768 & 12$\times$ & $2^{15}$ \\
BLS12-377 & 14 & 16{,}384  & 30$\times$  & $2^{14}$ \\
\bottomrule
\end{tabular}
\caption{Pathological (windowBits, curve) combinations on M3~Pro.
  Regression factor is the ratio of observed latency to the best adjacent
  window size at the same point count. Trigger column shows the approximate
  bucket count at which the pathology occurs.}
\label{tab:pathological}
\end{table}

\subsection{Root Cause Hypothesis}

The pathological bucket counts cluster near powers of two and their immediate
neighbors ($2^{10}{+}1$, $2^{11}{+}1$, $2^{14}$, $2^{14}{+}1$, $2^{15}$). We
hypothesize that the M-series tile-based GPU scheduler exhibits occupancy cliffs
when the threadgroup grid size aligns with boundaries of the GPU's internal
execution unit geometry.

Apple's M-series GPU is a tile-based deferred renderer~\cite{apple-gpu-arch}
with a compute dispatch model that maps threadgroups to execution units via an
internal scheduler whose behavior is not publicly documented. Unlike NVIDIA's
warp scheduler, which provides occupancy calculators and exposes register file
limits per SM~\cite{cuda-occupancy}, Metal provides no equivalent tooling or
documentation for compute workloads.

Register pressure amplifies the effect. BLS12-377's 12-limb field arithmetic
requires ${\sim}14$ temporary 32-bit values per Montgomery multiply, consuming
substantially more of the GPU register file per thread than BN254's 8-limb
arithmetic. This explains why BLS12-377 exhibits the most severe regression
(30$\times$) at the same bucket count where BN254 shows 140$\times$ but from a
higher baseline---fewer threads fit per execution unit, so the scheduling cliff
has a larger absolute impact.

We confirmed that this phenomenon is \emph{not} reproducible on NVIDIA GPUs.
ICICLE-CUDA~\cite{icicle} running on an RTX~3090~Ti shows smooth performance
scaling across the same bucket count range, consistent with CUDA's well-characterized
warp scheduler behavior.

\subsection{Mitigation}

We employ three mitigation strategies:

\paragraph{Per-curve window selection tables.}
Each curve maintains an empirically derived lookup table mapping point count
ranges to safe window widths that avoid pathological bucket counts. For example,
BN254 uses $w{=}16$ for all sizes above $2^{15}$; secp256k1 uses $w{=}13$ for
$\leq 2^{16}$ points and $w{=}16$ for larger inputs; BLS12-377 avoids $w{=}14$
entirely.

\paragraph{Signed-digit decomposition.}
Converting scalar digits to a signed representation $[{-}(2^{w-1}{-}1),\;
2^{w-1}]$ via a carry chain halves the effective bucket count, sidestepping some
pathological bands. This optimization was originally motivated by reducing
computation (42\% MSM speedup, \S\ref{sec:msm-opt}) but has the secondary
benefit of shifting the bucket count away from power-of-two boundaries.

\paragraph{Selective GLV disabling.}
GLV is only enabled when the resulting bucket count falls in a safe range. On
BN254, GLV's $2\times$ point increase with $w{=}16$ yields 32{,}768 signed
buckets, which is safe. On secp256k1, the same doubling hits pathological bands,
so GLV is disabled by default.

% Figure 3 placeholder
\begin{figure}[t]
\centering
\fbox{\parbox{0.95\columnwidth}{\centering\vspace{2em}%
\textbf{Figure 3:} Heatmap of MSM latency (log scale) as a function of
\texttt{windowBits} ($x$-axis) and point count ($y$-axis) for BN254,
secp256k1, and BLS12-377. Pathological bands appear as bright horizontal
stripes at specific bucket counts near powers of two. Safe operating regions
(dark) surround the cliffs.%
\vspace{2em}}}
\caption{MSM latency heatmap across window widths and point counts.
  Pathological scheduling bands are visible as 10--140$\times$ latency spikes
  at bucket counts near powers of two.}
\label{fig:heatmap}
\end{figure}

\paragraph{Recommendation.}
Apple should expose threadgroup scheduling hints or document the compute
occupancy model for M-series GPUs, as NVIDIA does for
CUDA~\cite{cuda-occupancy}. Absent this, developers of GPU-accelerated
cryptographic workloads must resort to exhaustive empirical sweeps to avoid
performance cliffs that are invisible to static analysis.

%% --------------------------------------------------------------------------
\section{Optimization Methodology}
\label{sec:optimization}

We present detailed case studies of three optimization efforts that together
account for the majority of zkMetal's performance gains. Each follows a
measure-first methodology: profile the hot path, form a hypothesis, implement
the optimization, and measure the result---discarding changes that do not
produce statistically significant improvement.

\subsection{MSM Optimization}
\label{sec:msm-opt}

Multi-scalar multiplication (MSM) is the dominant cost in KZG commitment
schemes, Groth16 proving, and IPA-based protocols. We optimized BN254~G1 MSM
from a 270\,ms baseline (naive GPU Pippenger at $2^{18}$ points on M3~Pro) to
45\,ms---a 6$\times$ improvement. Table~\ref{tab:msm-waterfall} and
Figure~\ref{fig:msm-waterfall} show the cumulative progression.

\begin{table}[t]
\centering
\small
\begin{tabular}{lrrr}
\toprule
\textbf{Optimization} & \textbf{Time} & \textbf{$\Delta$} & \textbf{Cum.\ Red.} \\
\midrule
Baseline (naive GPU Pippenger)  & 270\,ms &       & --- \\
Single command buffer           & 235\,ms & $-$13\% & 13\% \\
\texttt{fp\_sqr} removal        & 200\,ms & $-$15\% & 26\% \\
Segment tuning (512$\to$256)    & 173\,ms & $-$14\% & 36\% \\
Signed-digit decomposition      & 100\,ms & $-$42\% & 63\% \\
GLV endomorphism                &  78\,ms & $-$22\% & 71\% \\
CIOS \texttt{fp\_sqr} + Z=1 mixed add & 55\,ms & $-$30\% & 80\% \\
Register reuse + simplified ops &  45\,ms & $-$18\% & 83\% \\
Buffer caching                  &  45\,ms & $-$5\%  & 83\% \\
\bottomrule
\end{tabular}
\caption{MSM optimization waterfall for BN254 G1 at $2^{18}$ points (M3~Pro).
  Percentages are relative to the immediately preceding step. Cumulative
  reduction is relative to the 270\,ms baseline.}
\label{tab:msm-waterfall}
\end{table}

% Figure 4 placeholder
\begin{figure}[t]
\centering
\fbox{\parbox{0.95\columnwidth}{\centering\vspace{2em}%
\textbf{Figure 4:} Waterfall chart of MSM optimization progression from
270\,ms to 45\,ms with each optimization's contribution shown as a
cascading bar. Signed-digit decomposition and CIOS field arithmetic
contribute the two largest single improvements.%
\vspace{2em}}}
\caption{MSM optimization waterfall from 270\,ms to 45\,ms (BN254 G1,
  $2^{18}$ points, M3~Pro).}
\label{fig:msm-waterfall}
\end{figure}

\paragraph{Signed-digit decomposition.}
The standard Pippenger algorithm encodes each scalar as unsigned $w$-bit digits,
yielding $2^w - 1$ buckets per window. Signed-digit recoding converts each digit
to the range $[{-}(2^{w-1}{-}1),\; 2^{w-1}]$ via a carry chain, halving the
bucket count to $2^{w-1}$. The sign bit is encoded in bit~31 of the sorted
bucket index; the GPU reduce kernel checks this bit and negates the point's
$y$-coordinate via \texttt{fp\_neg} when set. This single optimization reduces
MSM latency by 42\% because the bucket accumulation phase---which is
memory-bandwidth-bound due to random-access scatter writes---processes half as
many buckets.

\paragraph{GLV endomorphism.}
BN254's efficiently computable endomorphism $\phi: (x, y) \mapsto (\beta x, y)$
with $\phi(P) = \lambda P$ allows decomposing a 256-bit scalar multiplication
into two 128-bit half-width multiplications~\cite{gallant-lambert-vanstone}. The
CPU decomposes each scalar via Babai rounding; a GPU kernel applies $\phi$ to
each point. This doubles the effective point count but halves the number of
windows from 16 to 8, yielding a net 22\% improvement. Critically, GLV is only
beneficial when point addition cost is low relative to the scalar width
reduction---it \emph{regresses} on BLS12-377 (1{,}155\,ms vs.\ 911\,ms at
$2^{18}$) because 12-limb point additions are so expensive that doubling the
point count costs more than halving the scalar width saves.

\paragraph{CIOS Montgomery and curve arithmetic.}
The original implementation used a Sum-of-Squares (SOS) squaring specialization
for \texttt{fp\_sqr}. While SOS saves ${\sim}25\%$ multiply operations for
squaring on CPUs, on the M3~GPU it requires 17 temporary registers versus CIOS's
10, degrading occupancy. Replacing \texttt{fp\_sqr} with standard CIOS
multiplication (\texttt{fp\_mul(a,a)}) improved MSM by 15\%. Additionally, we
introduced a specialized \texttt{point\_add\_mixed\_z1} routine that exploits
the fact that the first point added to each bucket has $Z = 1$ (affine),
saving one squaring and three multiplications per bucket (${\sim}262$K buckets at
$w{=}16$). Register reuse across point operation temporaries and a simplified
\texttt{fp\_neg} (single conditional subtraction instead of three) contributed
further gains.

\paragraph{Metal compiler \texttt{noinline} requirement.}
Functions returning large structs (\texttt{PointProjective} contains
three \texttt{Fp} fields totaling 96~bytes) are silently miscompiled when the
Metal shader compiler inlines them. We discovered this when 5 of 9 MSM
correctness tests failed after removing \texttt{\_\_attribute\_\_((noinline))}
annotations. The \texttt{noinline} attribute is mandatory on
\texttt{point\_double}, \texttt{point\_add\_mixed}, and \texttt{point\_add}
across all curves. This appears to be a genuine Metal compiler bug---the
generated code produces numerically incorrect results, not merely suboptimal
performance.

\paragraph{Timing breakdown at 45\,ms.}
The final timing breakdown for $2^{18}$ BN254 points is: GLV decomposition
2\,ms, precomputation 1.6\,ms, radix sort 2.4\,ms, GPU reduce (bucket
accumulation) 15.5\,ms, bucket sum + window combine 8.8\,ms, Horner evaluation
0.5\,ms. The GPU reduce phase consumes 50\% of total latency and is
fundamentally bottlenecked by random-access memory writes during the scatter
phase of bucket accumulation. This is consistent with the $9\times$ theoretical
headroom (\S\ref{sec:theoretical-floor}): the floor is set by memory bandwidth
for scatter writes, not compute.

\paragraph{Comparison.}
At $2^{18}$, zkMetal achieves 45\,ms versus Arkworks~CPU at 266\,ms
($5.9\times$), ICICLE-Metal at 1{,}475\,ms ($33\times$), and ICICLE~CPU at
556\,ms ($12\times$). ICICLE-CUDA on an RTX~3090~Ti achieves ${\sim}9$\,ms at
$2^{16}$~\cite{icicle}, a $3\times$ advantage attributable to native 64-bit
integer multiply and $6\times$ higher memory bandwidth (900~GB/s vs.\
150~GB/s).

\subsection{NTT Optimization}
\label{sec:ntt-opt}

The Number Theoretic Transform (NTT) is the core operation in polynomial
arithmetic for SNARKs and STARKs. We describe three GPU-level optimizations
and quantify the impact of field element size on throughput.

\paragraph{Four-step FFT.}
For sizes $\geq 2^{20}$, a single-kernel NTT exceeds Metal's 32\,KB
threadgroup shared memory limit. We adopt the four-step FFT
algorithm~\cite{bailey-ffts}: split an $N$-point NTT into row NTTs of size
$N_1$, a twiddle factor multiplication, and column NTTs of size $N_2$, where
$N = N_1 \times N_2$. Each sub-NTT fits in shared memory. The split point is
chosen so that $N_1 = N_2 = \sqrt{N}$ when possible, or the nearest power of
two otherwise.

\paragraph{Fused bitrev+butterfly.}
The standard NTT pipeline consists of a bit-reversal permutation followed by
$\log_2 N$ butterfly stages. We fuse the bit-reversal into the first butterfly
stage: the kernel reads input elements at bit-reversed indices from an input
buffer and writes butterfly results to a scratch buffer, eliminating one full
array traversal. This requires separate input and output buffers (in-place
fusion causes cross-threadgroup data races). The result is a 47\% improvement
for BN254 NTT at $2^{16}$ (0.85\,ms $\to$ 0.45\,ms) and 32\% at $2^{18}$
(2.44\,ms $\to$ 1.67\,ms). Gains diminish at larger sizes where global
butterfly stages dominate.

\paragraph{Twiddle fusion.}
In the four-step algorithm, the twiddle factor multiplication between row and
column NTTs is a separate dispatch over the full array. We fuse this
multiplication into the row NTT kernel: twiddle factors are applied during the
shared memory load phase (forward NTT) or during writeback (inverse NTT). This
eliminates a full-array memory pass of 64--128\,MB at large sizes. The impact is
dramatic at $2^{24}$: BabyBear NTT improves from 11.25\,ms to 2.32\,ms
($4.8\times$), and Goldilocks from 13.73\,ms to 3.00\,ms ($4.6\times$). BN254
benefits less because the row NTT kernel is already compute-bound at 32~bytes
per element.

\paragraph{Small-field throughput advantage.}
BabyBear (31-bit prime, $p = 2^{31} - 2^{27} + 1$) achieves 8.5 billion
elements per second at $2^{24}$, compared to 144~million for BN254---a
$59\times$ throughput ratio for an $8\times$ element size ratio. This
super-linear scaling arises because: (1) each BabyBear butterfly is a single
32-bit multiply-accumulate (MAD) versus 64~MADs for BN254 CIOS Montgomery; (2)
4-byte elements pack 8$\times$ more densely in cache lines, improving strided
column access utilization from 25\% (BN254, 32~bytes / 128-byte cache line) to
100\% (BabyBear, 4~bytes $\times$ 32 per line); and (3) BabyBear NTT at $2^{24}$
is within 1.2$\times$ of the memory bandwidth floor, indicating that computation
is no longer the bottleneck.

% Figure 5 placeholder
\begin{figure}[t]
\centering
\fbox{\parbox{0.95\columnwidth}{\centering\vspace{2em}%
\textbf{Figure 5:} NTT throughput (elements/sec, log scale) versus field
element size (bits) at $2^{24}$ elements. Points: BabyBear (31-bit,
8.5\,B\,elem/s), Goldilocks (64-bit, 5.7\,B\,elem/s), BN254 Fr (256-bit,
144\,M\,elem/s). The relationship is super-linear due to cache-line packing
and single-instruction arithmetic for small fields.%
\vspace{2em}}}
\caption{NTT throughput versus field element size at $2^{24}$.
  The 32-bit ALU is an \emph{advantage} for small fields.}
\label{fig:ntt-throughput}
\end{figure}

\paragraph{BN254 near-optimality at large sizes.}
BN254 NTT at $2^{22}$ runs in 26\,ms, which we believe is near-optimal for
32-byte elements on this hardware. The four-step algorithm's column NTTs access
memory at stride $N_1 \times 32$~bytes, yielding only 25\% cache-line
utilization. We tested three alternative approaches to improve coalescing:
transpose-before-column NTT (4 transposes + 4 blits cost ${\sim}50$\,ms,
exceeding the benefit), swapped row/column split dimensions, and individual
phase profiling. All three failed to improve over the baseline
(\S\ref{sec:dead-ends}).

\subsection{C CIOS CPU Acceleration}
\label{sec:ccios}

While GPU acceleration dominates for large inputs, many ZK primitives contain
CPU-side hot loops for small inputs, proof composition, or sequential operations
(Fiat-Shamir transcripts, scalar decomposition, polynomial evaluation). We
developed a systematic methodology for accelerating these paths.

\paragraph{Methodology.}
We profile each primitive to identify functions where $> 50\%$ of CPU time is
spent in field arithmetic. These functions are reimplemented in C using
\texttt{\_\_uint128\_t} for 64-bit CIOS Montgomery multiplication on ARM64.
The C compiler (\texttt{clang -O3}) maps \texttt{\_\_uint128\_t} multiplications
directly to ARM64 \texttt{MUL}+\texttt{UMULH} instruction pairs with optimal
register allocation and pipeline scheduling. For BN254's 256-bit field, each
Montgomery multiplication uses four 64-bit limbs with fully unrolled CIOS inner
loops, versus Swift's \texttt{multipliedFullWidth} + \texttt{addingReportingOverflow}
pattern which generates substantially more instructions due to overflow-reporting
semantics.

\paragraph{Key insight.}
Swift's optimizer handles 256-bit field arithmetic at 80--90\% of C performance
for \emph{isolated} multiplications. However, it is catastrophically slower (up
to $38\times$) for \emph{inner-loop} scalar operations where function call
overhead, overflow checking, and branch misprediction compound across millions of
iterations. The worst cases are tight loops: Pippenger bucket accumulation, NTT
butterfly inner loops, and GKR sumcheck rounds.

\begin{table}[t]
\centering
\small
\begin{tabular}{lrrr}
\toprule
\textbf{Primitive} & \textbf{Before} & \textbf{After} & \textbf{Speedup} \\
\midrule
BN254 NTT $2^{20}$  (CPU)    & 7.5\,s    & 247\,ms   & 30$\times$ \\
BN254 NTT $2^{14}$  (CPU)    & 103\,ms   & 2.7\,ms   & 38$\times$ \\
GKR $2^{10}$ d=4             & 241\,ms   & 7.7\,ms   & 31$\times$ \\
Verkle proof 256             & 670\,ms   & 5\,ms     & 134$\times$ \\
ECDSA batch 64 (CPU)         & 97\,ms    & 1.7\,ms   & 57$\times$ \\
HyperNova fold               & 3.6\,ms   & 0.09\,ms  & 40$\times$ \\
Spartan prove $2^{14}$       & 1{,}051\,ms & 121\,ms & 8.6$\times$ \\
Lasso prove $2^{18}$         & 481\,ms   & 56\,ms    & 8.2$\times$ \\
Lasso verify $2^{18}$        & 1.6\,s    & 31\,ms    & 52$\times$ \\
Groth16 prove 256            & 1.5\,s    & 14\,ms    & 107$\times$ \\
\midrule
BabyBear NTT $2^{22}$ (NEON) & 192\,ms   & 33\,ms    & 5.8$\times$ \\
Goldilocks NTT $2^{20}$ (C)  & 53\,ms    & 25\,ms    & 2.1$\times$ \\
\bottomrule
\end{tabular}
\caption{C CIOS acceleration results across zkMetal primitives. Top section:
  \texttt{\_\_uint128\_t} CIOS Montgomery for 256-bit fields. Bottom section:
  field-specific C optimizations (NEON SIMD for BabyBear, \texttt{\_\_uint128\_t}
  pipelining for Goldilocks). ``Before'' times use vanilla Swift field
  arithmetic.}
\label{tab:ccios}
\end{table}

\paragraph{Results.}
Table~\ref{tab:ccios} summarizes the impact across all primitives where C~CIOS
was applied. Speedups range from 2.1$\times$ (Goldilocks, where Swift's 64-bit
arithmetic is already efficient) to 134$\times$ (Verkle trees, where IPA's
iterated scalar multiplications compound the per-operation advantage across
$\log_2(n)$ rounds). The geometric mean speedup across all 256-bit primitives
is $32\times$.

Notably, C~Pippenger MSM on CPU achieves 68\,ms at $2^{16}$ BN254 points,
matching Arkworks~\cite{arkworks} (69\,ms) on identical hardware. This
demonstrates that the C~CIOS methodology alone---without any GPU
involvement---achieves state-of-the-art CPU performance for 256-bit field
arithmetic on ARM64.

\paragraph{NEON BabyBear NTT.}
For the 31-bit BabyBear field, we adapt the Plonky3 Montgomery
technique~\cite{plonky3} to ARM64 NEON: a 4-wide \texttt{uint32x4\_t}
Montgomery butterfly processes four field elements per instruction. The key
implementation subtlety is that the Montgomery constant must be the
\emph{positive} inverse $p^{-1} \bmod 2^{32}$ (not $-p^{-1}$), because the
NEON subtraction variant uses \texttt{vhsubq\_s32} which requires the
$(a \cdot b - q \cdot p)/2^{32}$ formulation. An initially incorrect sign on
\texttt{P\_INV} produced silently wrong results. The corrected NEON NTT achieves
5.8$\times$ over Swift, reducing $2^{22}$ BabyBear NTT from 192\,ms to 33\,ms.

\paragraph{Diminishing returns for 64-bit fields.}
Goldilocks ($p = 2^{64} - 2^{32} + 1$) gains only 2.1$\times$ from C
optimization because Swift already handles single-\texttt{UInt64} arithmetic
efficiently, and Goldilocks' special reduction (exploiting the sparse prime
structure) avoids the multi-limb carry chains where Swift's overhead is worst.
This confirms that C~CIOS acceleration is most impactful for 256-bit fields
where the 4-limb$\times$4-limb CIOS inner loop is the dominant cost.

\paragraph{Integration.}
All C code is compiled as a Swift Package Manager C target
(\texttt{NeonFieldOps}) with \texttt{-O3} and linked at build time. No
runtime JIT or assembly is involved. We verified that hand-written AArch64
assembly for Montgomery multiplication is 23\% \emph{slower} than
\texttt{clang~-O3} on \texttt{\_\_uint128\_t} code---the compiler already emits
optimal \texttt{MUL}+\texttt{UMULH} pairs. Developers should not hand-optimize
ARM64 Montgomery arithmetic.
